% Part: computability
% Chapter: computability-theory
% Section: reducibility

\documentclass[../../../include/open-logic-section]{subfiles}

\begin{document}

\olfileid{cmp}{thy}{red}
\olsection{Reduksibilitas}

\begin{explain}
Saiki kita ngerti yen ana saora-orane siji
himpunan, $K_0$, kang bisa dienumerasi
sacara komputabel nanging ora komputabel.
Mesthine ana himpunan liyane. Cara reduksibilitas
menehi sarana kang kuwat kanggo nuduhake
yen himpunan liyane nduwé sipat kaya mangkono,
tanpa kudu tansah bali menyang prinsip dhasar.

Umume, ``reduksi'' himpunan $A$ menyang
himpunan~$B$ yaiku cara ngowahi wangsulan
apa !!{element}s kalebu~$B$ dadi wangsulan
apa !!{element}s kalebu~$A$. Kita bakal
nggatekake jinis kang diarani ``reduksibilitas
akeh marang siji'', nanging isih ana
akeh jinis reduksibilitas liyane kanthi
sipat kang beda-beda. Pangerten reduksibilitas
uga wigati ing kajian kompleksitas
komputasional, kang kudu nggatekake
efisiensi. Umpamane, sawijining himpunan
diarani ``NP-komplit'' yen kalebu NP lan
saben masalah NP bisa direduksi menyang
himpunan iku, nganggo jinis reduksi
kang padha karo kang bakal dijlentrehake
ing ngisor iki, nanging kanthi syarat
tambahan yen reduksine bisa diitung
ing wektu polinomial.

Kita sejatine wis nggunakake pangerten
reduksi kanthi ora langsung. Tetepna
himpunan~$K$ minangka
\[
K = \Setabs{x}{\cfind{x}(x) \fdefined},
\]
yaiku $K = \Setabs{x}{x \in W_x}$.
Bukti yen masalah mandheg ora bisa
diputusake (\olref[hlt]{thm:halting-problem})
kanthi paling langsung nuduhake yen
$K$ ora komputabel. Elinga yen $K_0$
minangka himpunan
\[
K_0 = \Setabs{\tuple{e, x}}{\cfind{e}(x) \fdefined },
\]
% OLPL-141: Sumber beku mbalikke urutan pasangan ing katrangan W_e; definisi K_0 ing ndhuwur lan teges x kalebu W_e mbutuhake (e,x).
yaiku $K_0 = \Setabs{\tuple{e,x}}{x \in W_e}$.
Gampang ngembangake bukti yen~$K$
ora komputabel dadi bukti yen~$K_0$
ora komputabel: yen $K_0$ komputabel,
kita bisa mutusake apa !!a{element}~$x$
kalebu $K$ mung kanthi takon apa
pasangan $\tuple{x, x}$ kalebu $K_0$.
Fungsi~$f$ kang ngirim $x$ menyang
$\tuple{x, x}$ minangka conto
\emph{reduksi} saka $K$ menyang $K_0$.
\end{explain}

\begin{defn}
Tetepna $A$ lan $B$ minangka himpunan
wilangan asli. Fungsi komputabel
~$f\colon \Nat \to \Nat$ diarani
\emph{reduksi akeh marang siji} saka $A$
menyang~$B$ yen lan mung yen, kanggo
saben wilangan asli~$x$,
\[
x \in A \quad \text{yen lan mung yen} \quad f(x) \in B.
\]
Yen ana reduksi $f$ kaya mangkono,
kita nyebut $A$ \emph{bisa direduksi akeh
marang siji} menyang~$B$, ditulis
$A \leq_m B$. Yen $A$ bisa direduksi
akeh marang siji menyang $B$ lan
kosok baline, $A$ lan $B$ diarani
\emph{setara miturut reduksi akeh marang siji},
ditulis $A \equiv_m B$.
\end{defn}

Yen fungsi $f$ ing definisi ndhuwur
jebul injektif, $A$~diarani
\emph{bisa direduksi siji marang siji}
menyang~$B$. Umume reduksi kang
bakal diterangake ing ngisor iki
nyukupi syarat kang luwih kuwat iku,
nanging kita ora bakal nggunakake
kasunyatan kasebut.

\begin{digress}
Pancen bener, nanging ora cetha,
yen reduksibilitas siji marang siji
iku syarat kang tenan luwih kuwat
tinimbang reduksibilitas akeh marang siji.
Kanthi tembung liya, ana himpunan
tanpa wates $A$ lan~$B$ kang ndadekake
$A$ bisa direduksi akeh marang siji
menyang~$B$, nanging ora bisa direduksi
siji marang siji menyang~$B$.
\end{digress}

\end{document}
